\documentclass[12pt]{article}
\usepackage[utf8]{inputenc}
\usepackage[T1]{fontenc}
\usepackage{hyperref}
\title{Northern Uzbek Language Resources for K-12 Curriculum Development}
\author{Compiled by Claude Sonnet 5 (claude-sonnet-5) for Andr\'as Kornai}
\date{}
\begin{document}
\maketitle

\section{Introduction}
Northern Uzbek (uzn; Glottocode nort2690) is a Karluk Turkic language and the official language of Uzbekistan, used as a language of instruction. It is also spoken in neighbouring countries. This dossier is limited by the search tool, which returned mostly English-language and foreign-learner material. I could not confirm current titles or grade coverage of the national school textbooks for the Uzbek-language subject (Ona tili), and a check of the ministry's own site (uzedu.uz) is still needed.

\section{Primary Educational Resources}
\begin{itemize}
\item RTI International, Desk Review of Uzbek Language Arts and Mathematics Student Textbooks and Teacher Guides, covering the grades 1--4 pilot textbooks produced with Uzbekistan's Ministry of Preschool and School Education under the USAID program: \url{https://shared.rti.org/node/1252}
\item RTI International, Textbook and Teacher Guide Development, a collaborative approach: \url{https://shared.rti.org/node/1246}
\item Kun.uz news, USA to provide \$10 million worth of textbooks to Uzbekistan (2021): \url{https://kun.uz/en/news/2021/08/07/usa-to-provide-10-million-worth-of-textbooks-to-uzbekistan}
\item Uzbekistan government open data portal, dataset 16487 (content not verified): \url{https://olddata.gov.uz/en/datasets/16487}
\end{itemize}

\section{Linguistic and Reference Materials}
\begin{itemize}
\item Wikipedia, Uzbek language: \url{https://en.wikipedia.org/wiki/Uzbek_language}
\item University of Chicago CEERES, Uzbek Language Resources (lists an Indiana University online Uzbek-English dictionary and the Uzbek-Uzbek dictionary Savodxon Lu'gati): \url{https://ceeres.uchicago.edu/uzbek-language-resources}
\item Georgetown University Press, Uzbek: An Elementary Textbook, front matter (for foreign learners): \url{https://press.georgetown.edu/files/materials/9781589017061_Frontmatter.pdf}
\end{itemize}

\section{Religious and Community Resources}
None found in this search.

\section{Digital and Audio Resources}
\begin{itemize}
\item Omniglot Uzbek book listing: \url{https://www.omniglot.com/books/language/uzbek.htm}
\item Dunyo.info, report on Otabek Alimov's self-study guide in 22 lessons, which won a national competition (2025): \url{https://dunyo.info/en/undefined/%D0%8Ezbekiston-diplomati-tomonidan-%D1%9Ezbek-tilini-%D1%9Ergan-}
\end{itemize}
No Uzbek corpus was found in search results.

\section{Curriculum Development Context}
Uzbek is the language of instruction, so the main task is a curriculum built on existing national standards and textbooks rather than a revival. The USAID and RTI project documents describe standards-based Uzbek Language Arts materials for grades 1--4 and the handover of Native Language standards. Materials for grades 5--11 were not found. Both Latin and Cyrillic scripts are in use in different settings.
\end{document}